% \section{Modeling RCS Data with a Univariate Time Series}


There also exist aggregate methods that involve calculating 
averages of independent and dependent variables at each time 
step and analyze these averages as a time series with covariates.
 Such a model is an example of a state-space model that can be 
 modeled jointly via Kalman filtering \cite{BrockwellDavis1991}. 
 In the regression setting, aggregates would be needed for both 
 response and explanatory variables, fitting the following model:

$$ \bar{\mathbf{y}} = (n^{-1}X_t'X_f)\beta + \mathbf{u}, \qquad \mathbf{u}\sim\text{ARMA(p,q)} $$

  While this method is effective at estimating population-level 
  fixed effects like intercepts and slopes, but is limited when 
  it comes to the estimation of individual-level effects, 
  especially if indicator variables vary little in their 
  distribution at each time step. This is a highly plausible
   scenario when it comes to sampling methods for surveys, as 
   researchers interested in the effect of a certain covariate 
   may very well aim to be consistent in the proportions of time 
   step samples that have this feature.

  % $$ y_{ti} = \beta_0 + \beta_1t + \beta_2x_{ti} + u_t + \varepsilon_{ti}, \qquad i = 1,...,n_t$$

  % $$ u_t - \phi u_{t-1} = \eta_t + \theta \eta_{t-1}, \qquad t = 1,...,T$$ 
  
  
  % $$\eta_t \overset{\text{i.i.d.}}{\sim}N(0,\sigma_\eta^2), \qquad x_{ti} \sim \text{Bernoulli}(.5), \qquad \varepsilon_{ti} \overset{\text{i.i.d}}{\sim} N(0,\sigma_\varepsilon^2)$$






